\section{User Authorization}
\label{sec:authorization}

Zcash shielded addresses lack a standard message-signing mechanism. Even if a signature scheme existed, a raw signature does not prove temporal liveness or prevent replay attacks without a persistent, on-chain nonce registry.

ZNS solves this with an in-band one-time passcode (OTP) that relies exclusively on Zcash's native encryption. The Mint encrypts a 16-byte random payload to the Unified Address currently on record. Returning that exact payload to the Mint proves two things simultaneously: cryptographic read-access to the address, and temporal liveness bounded by the OTP expiration block.

A \action{claim} establishes initial state and requires no OTP. An \action{update} or \action{release} requires a valid OTP exchange before the Mint creates the Name Note \cite{zns-mint}.

\begin{center}
\znsauthorizationflow
\end{center}
\subsection{Authorization target}

Each action uses the following authorization address:

\begin{center}
\begin{tabularx}{\textwidth}{@{}lYY@{}}
\toprule
Action & Authorization address & Consequence \\
\midrule
\action{claim} & Proposed Unified Address & The request itself authorizes the claim; no OTP is issued. \\
\action{update} & Unified Address in the accepted live tip & The OTP tests access to the address that currently controls the name. \\
\action{release} & Unified Address in the accepted live tip & The OTP tests access to the address that currently controls the name. \\
\bottomrule
\end{tabularx}
\end{center}

The active manifest MUST identify the shielded receiver type used for OTP delivery.
The Mint selects that receiver from the controller Unified Address under ZIP~316 and rejects an address that does not contain it \cite{zip316}.
A successful OTP response establishes that the requester can receive shielded memos sent to that receiver.
It does not establish spending authority, control of every receiver in the Unified Address, or control of an unverified update destination.

The proposed new address in an \action{update} is not verified by the OTP.
A mistyped or attacker-substituted destination can become the live binding if wallet confirmation fails.

The public Name Note account MUST NOT carry OTPs because its full viewing key is public.
The Mint sends OTP relay memos from an account whose viewing authority is not published.
This separation prevents disclosure through the Name Note viewing key.
It does not hide network metadata from the Mint operator, wallet provider, or transaction-submission infrastructure.

\subsection{OTP relay memo}

The Mint sends the OTP to the selected receiver of the controller Unified Address in an Orchard shielded memo.
The relay memo grammar is:

\begin{center}
\begin{tabularx}{\textwidth}{@{}lY@{}}
\toprule
Non-padding bytes \\
\midrule
\field{ZNS:otp:<name>:<verb>:<ua>:<otp>} \\
\bottomrule
\end{tabularx}
\end{center}

The \field{<verb>} field is \field{update} or \field{release}.
The \field{<ua>} field is the target Unified Address from the request.
The \field{<otp>} field is exactly 32 lowercase hexadecimal characters encoding 16 bytes.

The relay memo is exactly 512 bytes and zero-padded.
A request parser MUST reject a relay memo as a request memo because its verb is \field{otp}.

\subsection{OTP properties}

An OTP is 16 bytes of entropy drawn from a cryptographically secure random source inside the attested environment \cite{zns-mint}.
The active manifest MUST fix the validity window \(N\) in blocks; the current implementation uses \(N = 24\).
An OTP issued at block \(h\) is valid through block \(h + N\).

The Mint verifies a provided OTP with constant-time equality, accepts it only for the exact \((\field{name}, \field{action}, \field{ua})\) it was issued for, and burns it on acceptance.
A consumed OTP MUST never succeed again.
The Mint MUST reject an expired OTP and an OTP for a mismatched transition.

The active manifest MUST fix the validity window, the persistent-state format, and the maximum number of attempts.
These values are \openstatus{} until the manifest fixes them.

\subsection{Request binding}

An unbound OTP could authorize a substituted request.
The Mint binds each OTP to the exact \((\field{name}, \field{action}, \field{ua})\) tuple in the first request.
It accepts the OTP only in a second request with the same tuple.

The Mint MAY impose additional request-binding rules such as a minimum request value or a payment locator.
Those rules are Mint policy, not independently verifiable by a Resolver.

\subsection{Public verification boundary}

The OTP exchange does not produce public per-request authorization evidence.
A Resolver can verify the resulting Name Note, Mint self-send, and Mint attestation.
It cannot reconstruct the OTP exchange from the public Name Note log.

Attestation establishes that a recognized Mint identity runs approved policy under the stated TEE assumptions.
It does not disclose a transcript proving that one OTP was delivered and consumed.
Per-request public authorization evidence requires a separate proof construction.

